\section{Spline smoothing of vector measurements}
\label{sec:vector}

Consider the {\em vector} measurement model:
\be
\y_n = \mfun(\g(t_n)) + \bveps_n,
\ \ n=1, \ldots, N
,\ee{eq:multi-model}
\[
\g(t_n) \in \Re^M, \ \ 
\bveps_n, \y_n \in \Re^L, \ \ 
\bveps_n \sim N(\bzero, \mcov_n), \ \
E \{ \bveps_n \bveps_m' \} = 0, \ n \neq m,
\]
with known symmetric and positive-definite
error covariances $\mcov_n$.
The goal of smoothing is to estimate 
the smooth vector-valued function
\g,
given the measurements $\{\y_n\}_{n=1}^N$.

The nonparametric approach to this problem
prescribes a compromise
between
the conflicting goals of
fit to the data
and
smoothness of the estimated functions.
An estimate is obtained
by using
the following criterion,
which is the natural generalization
of a similar criterion used for scalar measurements:
\be
\bghata = 
\arg \min_{\g}
\sum_{n=1}^N (\y_n - \mfun(\g(t_n)))'
\mcov_n^{-1} (\y_n - \mfun(\g(t_n)))
+ R_{\balpha}(\g)
,\ee{eq:multi-min}
where
\[
R_{\balpha}(\g) = \sum_{m=1}^M \alpha_m \int (\ddot{g}_m(t))^2 \ dt .
\]
The smoothing parameter $\balpha$
controls the tradeoff
between residual error and smoothness.
Although
$\bghata(t)$
is a continuous function,
it is uniquely determined by its values at the knots
$t_1, \ldots, t_N$,
denoted
$\hat{\y}_n = \bghata(t_n)$.

The {\tt vspline} library
contains special purpose subroutines
for efficiently minimizing~(\ref{eq:multi-min}).
Two cases are implemented:
``linear'' smoothing for $\mfun(\g)=\g$,
and ``nonlinear'' smoothing
where the nonlinear function $\mfun$
and its derivatives must be implemented by the user.

Subroutine
{\tt vspline},
described below,
computes
$\{ \hat{\y}_n \}_{n=1}^N$
given
$\balpha$,
%$\alpha_1, \ldots, \alpha_M$,
$\{ \y_n \}_{n=1}^N $,
$\{ t_n \}_{n=1}^N $,
and
$\{ \mcov_n^{-1} \}_{n=1}^N $.
The computational requirements~\cite{Fessler:91:lin} are $O(M^3\cdot N)$.
